\chapter{Analitički okvir za evaluaciju informacionih sistema}
U poglavlju \ref{sec:uvod} informacioni sistemi su analizirani kroz arhitektonsku organizaciju, modele skaliranja i savremene pristupe izvođenju u oblaku. Bez obzira na to da li sistem funkcioniše kao monolit, troslojna aplikacija ili kao skup instanci raspoređenih u cloud okruženju, takvi sistemi i dalje predstavljaju \textit{logički centralizovane informacione sisteme} u kojima postoji jedinstveni skup komponenti koje zajednički realizuju poslovnu logiku \cite{van2017distributed}. Čak i kada se više instanci izvršava na različitim uređajima, sistem se ponaša kao jedna integrisana celina sa jedinstvenim stanjem i kontrolom.

Za sistematsku analizu prednosti i ograničenja ovakvih informacionih sistema potrebno je fokusirati se isključivo na njihove \textbf{tehničke atribute infrastrukture}, nezavisno od menadžerskih, organizacionih i ekonomskih aspekata. Ovi aspekti biće samo ukratko pomenuti radi potpunosti, ali ne predstavljaju predmet dublje razrade u ovoj disertaciji.

U oblasti softverskog inženjerstva kvalitativne karakteristike informacionih i softverskih arhitektura formalizuju se kroz tzv. \textbf{atribute kvaliteta} \cite{bass2021software}, koji predstavljaju standardni analitički okvir za ocenu tehničkog ponašanja i sposobnosti sistema. Među najznačajnijim atributima za evaluaciju informacionih sistema izdvajaju se:
\begin{enumerate}
\item performanse i efikasnost,
\item skalabilnost,
\item dostupnost i otpornost na greške,
\item bezbednost i integritet,
\item kompleksnost implementacije, održavanja i upravljanja.
\end{enumerate}

Ovi atributi obezbeđuju tehnički konzistentan i teorijski utemeljen okvir za procenu kvaliteta savremenih informacionih sistema. Njihovim posmatranjem moguće je jasno razlikovati \textbf{tehnička svojstva sistema}, koja predstavljaju primarni predmet istraživanja u ovoj disertaciji, od šireg menadžerskog, organizacionog i ekonomskog konteksta, koji se ovde samo dopunski navodi.